% !TeX root = ../papers/nc-p3.tex
\section{非被約な配置で行列式が恒等零になるための局所長条件}
\label{sec:singular-opposite-root-lengths}
\begin{theorem}\label{thm:singular-opposite-root-lengths}
四辺有限な標準有限Frobenius窓で、最初と反対の配置は標識を保って射影非同値とする。
その対称二次式行列 $N$ の行列式が恒等零なら、最初の配置 $X$ の幾何点数は高々13である。
13点なら十二単純点と埋込次元3の一つの長さ8の点からなる。
また全局所環が曲線状なら、各局所長は偶数でなければならない。
その場合、各局所長を半分にした長さ10の部分スキームは、
三つの二次式で定義される次数3の決定式曲線に含まれる。
従ってこれらの必要条件を外れる配置でも、異なる反対位置は
$H^0(\mathcal I_X^3(8))$ に非零の八次式を与える。
\end{theorem}
\begin{proof}[核座標とその零点長]
一般階数は既出の排除により3であり、
$\operatorname{adj}N=hgg^{\mathsf t}$ と原始に書く。
$g$ の次数 $e$ は2または3である。線型の場合は混合双一次式が射影同定を与える。
$g$ の成分の張る次元 $r$ は3か4である。二次元以下の場合は
定理\ref{thm:opposite-sextic-threshold}の先頭二列の議論で排除できる。

まず $r=3$ とし、第四核成分を零にする。左上の対称三行を $H$、第四列の上三成分を $b$ とする。
$H$ の六成分は独立、一般階数は2以下である。階数1なら
$H=vv^{\mathsf t}$ で $v$ は独立な三つの一次式である。
$N|_X$ の階数1条件により、$v\times b$ の各三次成分は $X$ 上で各 $v_i$ に消される。
この消滅部分の次元は高々1である。実際完全跡でその双対は三つの一次式の共通点による
$X$ の切断の関数環である。三次評価が単射なので $v\times b$ は多項式としても
一定ベクトルの倍となる。$v$ との内積が零で三成分が独立だから恒等零である。
すると $b$ は $v$ の倍で $N$ の一般階数が2以下となり矛盾する。
従って $H$ は一般階数2で、$\operatorname{adj}H$ は原始な三成分 $g'$ の平方の多項式倍である。
次数4より $e=2$ でその倍数は非零定数。特に $G=(g'_1,g'_2,g'_3)$ は
$G^2\subset\mathcal I_X$ を満たす。

三つの独立で原始な二次式の一次シジジーの次元は高々2である。
分数体上独立な二つの一次シジジーがあれば、その外積は $g'$ の定数倍である
（原始性と次数）。どの第三の一次シジジーとも外積を取れば、分数体上の展開係数が定数となる。
分数体上の階数1なら、定数シジジーがないので原始ベクトルの定数倍に限られる。
従って $S_1G$ の次元は10以上。三次評価は $X$ で同型なので
$X\cap V(G)$ の長さは高々10である。特に台は高々十点である。
全点が曲線状なら各局所長 $n$ に対し $G^2|_X=0$ からこの切断の局所長は
$\lceil n/2\rceil$ 以上である。総和が10以下なので全て偶数で切断長は丁度半分である。
等号より二本の独立な一次シジジーがあり、先の外積で三つの二次式が
$3\times2$ 一次式行列の最大小行列式になる。原始性と行列式の高さ上界より高さ2で、
Hilbert--Burch分解から次数3の曲線となる。

$r=4,e=2$ の場合の $G$ を四成分の張る空間とする。このとき
\[
              \dim S_1G\ge12,\qquad \operatorname{length}(X\cap V(G))\le8 .
\]
実際一般平面への四二次式の制限は独立である。そうでなければ一般の平面方程式が
$G$ のある元を割る。各可約二次式の因子の選び方は有限なので、
$\mathbb P(G)$ の三次元全体が可約形式からなる。しかし可約二次式だけの線型空間は
固定因子を持つか固定二変数の二次式に入る。
これを見るには一つの階数2形式 $lm$ を固定する。別の階数2形式との鉛筆が階数2以下なら、
変数の二平面が同じか、一次因子を共有する。二平面が交わらなければ和の階数は4、
直線 $kx$ で交われば各形式を $(x,y),(x,z)$ の二変数に書ける。
三行列式の鉛筆パラメータの一次・二次係数から両者の $y^2,z^2$ 係数は零となる。
ゆえに線型空間は $l$ の倍元、$m$ の倍元、二変数形式の三つの線型空間の和集合に入り、
その一つに含まれる。全て階数1の場合も次元は高々1である。
四成分の独立性と原始性からどちらも不可能。
三変数で四つの相異なる二次先頭単項式の一次倍は少なくとも八種類である。
これに平面方程式と $G$ との積の四次元を足す。切断長の上界は再び三次評価から従う。

$h$ が平方でない二次式なら、$V(G)$ の外にある $X$ の台は $h$ の特異点に限られる。
三次小行列式 $hg_i^2$ は $X$ で二重に消えるからである。
特異軌跡は高々直線で、その $X$ 切断は長さ3以下、従って台は合計高々11点である。
さらにこの場合に曲線状 $X$ は存在しない。各長さ $n$ の局所点で形式座標を選んで
$I_X=(t^n,y,z)$ とし、$t,y,z$ に重み $1,n,n$ を与える。
$G$ 切断長を $j$、$h$ の特異軌跡による切断長を $d$ とする。
$j=n$ または $d=n$ なら自明に $2j\ge n-d$ である。
他の場合はある $g_i$ の重みが $j$、ある一次偏導 $\partial h$ の重みが $d$。
元の座標の偏微分は形式座標でも重みを高々 $n$ だけ下げるから $h$ の重みは $n+d$ 以下である。
$hg_i^2\in I_X^2$ により $2j\ge n-d$ を得る。
アフィンチャートで同次偏導を用いる場合もEuler式から、零階または一階の微分作用素でよい。
全点で足すと $\sum j\ge(20-3)/2$ となり上界8に反する。
\end{proof}

\begin{proof}[四本の三次式が整閉包に入る場合]
残るのは $r=4$ で $e=3$、または $e=2$ で $h$ が平方の場合である。
定数倍を吸収すれば余因子は $gg^{\mathsf t}$ または $(\ell g)(\ell g)^{\mathsf t}$。
いずれも四つの独立な三次式の各平方が $\mathcal I_X^2$ に入るので、
各式は局所イデアルの整閉包 $\overline{\mathcal I_X}$ に入る。
三次評価が同型だから
\[
  4\le\sum_{p\in X}\operatorname{length}(\overline{I_p}/I_p).
\]
曲線状イデアル $(t^n,y,z)$ は上の重みが $n$ 以上の元全体であり整閉であるから除かれる。

埋込次元2以下の長さ $n$ の局所完全交叉では
$\operatorname{length}(\overline{I_p}/I_p)\le\lfloor(n-1)/2\rfloor$ である。
一次変数を消去して二変数正則局所環に移る。整従属は商に移り、
残る二生成元イデアルとその整閉包 $J$ の重複度は同じ $n$。
二変数の有限余長イデアル $J$ では $e(J)<2\operatorname{length}(R/J)$ である。
実際局所単項式順序での初期イデアル $M$ への退化は余長を保ち、
$M^a\subseteq\operatorname{in}(J^a)$ より $e(J)\le e(M)$。
単項式の場合はNewton多角形の補の面積の二倍が重複度である。
その面積は標準単項式の単位正方形の個数より真に小さい。二軸切片を結ぶ線分の上方が
Newton多角形に入るため、最初の横軸切片の直前の単位正方形に正の差がある。
これで上界を得る。埋込次元3なら長さは少なくとも8である。
三つの局所生成元の位数は2以上で二次の随伴商は少なくとも三次元、
局所Gorenstein環のsocleは一次元だから三次以上の部分も非零である。

幾何点数が14以上なら $\sum(n-1)\le6$、従って全点が平面的で整閉包の増分は高々3。
13点でも全点が平面的なら同じ上界である。残る場合は余分な長さ7を全て
埋込次元3の一点で使う。最後に一般階数4なら、$X$ 上の単元成分を消去して残る三行の成分が
$\mathcal I_X$ に入ることから行列式の三重条件がスキームとして従う。
\end{proof}
